\subsection*{Wnioski}
\begin{itemize}
\item Transfer wiedzy z ImageNet daje wyraźnie lepsze wyniki niż sieć uczona od zera
  (macro-F1 w CV \cvFRes{} i \cvFEff{} wobec \cvFCnn{}). ResNet50 i EfficientNetB0 są
  praktycznie równorzędne, przy czym EfficientNetB0 jest ok.~6 razy mniejszy. Strojenie
  Optuną poprawiło wyniki nieznacznie, w granicach niepewności.
\item Najlepszy model (ResNet50 po strojeniu) osiąga na oryginalnym teście accuracy
  \accTestResTuned{}, a na teście oczyszczonym z duplikatów \accDedupResTuned{}
  (macro-F1 \fDedupResTuned{}). Różnica pokazuje, jak bardzo duplikaty między treningiem
  i testem zawyżają wyniki publikowane dla tego zbioru.
\item Model bardzo dobrze rozpoznaje guzy przysadki, oponiaki i obrazy bez guza, ale gubi
  glejaki pochodzące z innego źródła niż dane treningowe (czułość dla glejaka
  \sensDedupGlioma{}). Część z nich uznaje za obraz bez guza. Analiza błędów i Grad-CAM
  wskazują, że model częściowo nauczył się rozpoznawać \emph{źródło i sekwencję} obrazu
  zamiast cech guza.
\item W praktyce klinicznej taki model mógłby co najwyżej wspierać wstępne sortowanie badań
  (triage), z wynikiem zawsze weryfikowanym przez radiologa. W obecnej postaci nie nadaje się
  do użycia: przeoczenie co piątego glejaka z nowego źródła jest nieakceptowalne.
\end{itemize}

\subsection*{Problemy i ograniczenia}
\begin{itemize}
\item \textbf{Brak identyfikatora pacjenta.} Nie da się zrobić podziału per pacjent.
  Grupowanie niemal-duplikatów pHash jest tylko przybliżeniem: sąsiednie przekroje tego
  samego pacjenta mogą różnić się bardziej niż próg.
\item \textbf{Zbiór złożony z kilku źródeł}, w których klasy są nierówno reprezentowane
  (glejaki w treningu tylko z jednego źródła, obrazy bez guza głównie z innych). To główna
  przyczyna błędów. Do tego sztuczny balans (kopie \texttt{aug}) i szum etykiet
  (\nPairsConflict{} par niemal identycznych obrazów z różnymi etykietami).
\item \textbf{Pojedyncze przekroje 2D} zamiast całych badań 3D i różnych sekwencji tego samego
  pacjenta; radiolog ocenia całe badanie.
\item \textbf{Brak walidacji zewnętrznej} (inny szpital, inny skaner). Oczyszczony test pochodzi
  z tych samych źródeł co trening i jest mały (\nTestDedup{} obrazów, w tym tylko \nDedupNotumor{}
  bez guza), stąd szerokie przedziały ufności.
\item \textbf{Strojenie} na tych samych foldach co ocena CV (brak zagnieżdżonej walidacji).
\item Rozdzielczość 224$\times$224 może ukrywać drobne zmiany. Wyjaśnienia Grad-CAM mają niską
  rozdzielczość (7$\times$7) i pokazują korelacje, nie przyczyny.
\end{itemize}

\subsection*{Aspekty etyczne i prawne}
Obrazy medyczne to dane szczególnej kategorii (art.~9 RODO). Ich użycie wymaga podstawy
prawnej i anonimizacji, a pacjent ma prawo do informacji o zautomatyzowanym przetwarzaniu
i do wyjaśnienia decyzji (art.~22). Uprzedzenia ukryte w danych -- tu powiązanie klasy ze
źródłem obrazu -- nie są widoczne w samych metrykach i mogą prowadzić do systematycznie
gorszej diagnozy dla pacjentów badanych w innych ośrodkach. Odpowiedzialność za decyzję
zawsze ponosi lekarz.

\subsection*{Możliwości rozwoju}
\begin{itemize}
\item Dane: zbiory z identyfikatorem pacjenta i wieloma ośrodkami (np. BraTS -- badania 3D
  z wieloma sekwencjami), walidacja na zewnętrznym zbiorze, glejaki z różnych źródeł w treningu.
\item Walidacja grupowa według źródła (leave-one-source-out), żeby wprost mierzyć
  przenoszenie się modelu na nowy ośrodek.
\item Normalizacja intensywności per obraz (np. z-score w obrębie mózgu), usunięcie ramek
  i znaczników, wycinanie samego mózgu -- żeby utrudnić uczenie się skrótów.
\item Segmentacja guza (np. U-Net) zamiast samej klasyfikacji, modele 3D.
\item Progi decyzyjne dobrane pod wysoką czułość klas nowotworowych, kalibracja
  (temperature scaling) i wykrywanie obrazów spoza znanej dystrybucji (odmowa odpowiedzi
  przy niepewności).
\end{itemize}
